% Generated draft: requires mathematical and source review.
\documentclass[11pt]{article}
\usepackage[margin=25mm]{geometry}
\usepackage{fontspec}
\setmainfont{Latin Modern Roman}
\usepackage{amsmath,amssymb,mathtools,unicode-math}
\setmathfont{Latin Modern Math}
\usepackage{enumitem,needspace,xurl,hyperref}
\setlength{\parindent}{0pt}
\setlength{\parskip}{5pt}
\setlength{\emergencystretch}{3em}
\newcommand{\N}{\mathbb{N}}
\newcommand{\Z}{\mathbb{Z}}
\newcommand{\Q}{\mathbb{Q}}
\newcommand{\R}{\mathbb{R}}
\newcommand{\C}{\mathbb{C}}
\newcommand{\F}{\mathbb{F}}
\newcommand{\A}{\mathbb{A}}
\newcommand{\PP}{\mathbb{P}}
\newcommand{\E}{\mathbb{E}}
\newcommand{\eps}{\varepsilon}
\newcommand{\dd}{\,\mathrm{d}}
\begin{document}
\section*{UnsolvedMath ID: 30003228}
\textbf{Source problem code:} catalogue entry
\subsection{Definitions and mathematical statement}
Let \(F=(f_1,\ldots,f_r):(\C^n,0)\to(\C^r,0)\) be a holomorphic germ and put \(f=\prod_{i=1}^r f_i\). The supplied record does not explicitly exclude identically zero components. The following standard definitions require \(f\not\equiv0\); this necessary qualification must be checked against the source before accepting the formulation.

Write \(\mathbf s=(s_1,\ldots,s_r)\), \(\mathbf1=(1,\ldots,1)\), and \(F^{\mathbf s}=\prod_i f_i^{s_i}\), interpreted as a formal symbol over \(\mathcal O_{\C^n,0}[f^{-1}][\mathbf s]\). Let \(\mathcal D_{\C^n,0}\) denote the ring of germs of holomorphic differential operators. Using the simultaneous-shift convention, define the local Bernstein--Sato ideal by
\[
 B_F=\left\{b(\mathbf s)\in\C[\mathbf s]:
 b(\mathbf s)F^{\mathbf s}=P(\mathbf s)F^{\mathbf s+\mathbf1}
 \text{ for some }P(\mathbf s)\in\mathcal D_{\C^n,0}[\mathbf s]\right\}.
\]
Here differentiation acts through
\[
 \partial_{x_j}F^{\mathbf s}
 =\left(\sum_{i=1}^r s_i\frac{\partial_{x_j}f_i}{f_i}\right)F^{\mathbf s}.
\]
Set
\[
 Z(B_F)=\{\mathbf s\in\C^r:b(\mathbf s)=0\text{ for every }b\in B_F\},
 \qquad
 \operatorname{Exp}(\mathbf s)=(e^{2\pi i s_1},\ldots,e^{2\pi i s_r}).
\]
The supplied record does not explicitly specify the shift convention for \(B_F\).

For a sufficiently small representative ball \(B_\epsilon\) about the origin, let \(U_\epsilon=B_\epsilon\setminus f^{-1}(0)\). For \(\boldsymbol\lambda\in(\C^*)^r\), let \(\mathcal L_{\boldsymbol\lambda}\) be the rank-one complex local system on \((\C^*)^r\) whose monodromy around the \(i\)-th coordinate circle is \(\lambda_i\). A precise support locus for the complement at the origin is
\[
 V_{F,0}=\left\{\boldsymbol\lambda\in(\C^*)^r:
 H^q\!\left(U_\epsilon,(F|_{U_\epsilon})^*\mathcal L_{\boldsymbol\lambda}\right)\ne0
 \text{ for some integer }q\ge0\right\}.
\]
The source denotes a cohomology-support locus by \(V_F\), but the supplied material does not determine whether it means this locus at the origin or a union of analogous local loci over nearby points of \(f^{-1}(0)\). Accordingly, \(V_F\) below retains the source's unresolved convention.

The original source record presents the following three assertions for every germ satisfying the intended hypotheses:
\begin{enumerate}
\item There are finitely many generators \(g_1,\ldots,g_m\) of \(B_F\), each of which is a finite product of affine forms
\[
 a_1s_1+\cdots+a_rs_r+b,
 \qquad a_i\in\Q_{\ge0},\quad b\in\Q_{>0}.
\]
The supplied coefficient conditions do not require \((a_1,\ldots,a_r)\ne0\).
\item The equality
\[
 \operatorname{Exp}(Z(B_F))=V_F
\]
holds with the source's intended definition of \(V_F\).
\item Every polar hyperplane \(H\) of the multivariable motivic zeta function associated with \(F\) satisfies
\[
 H\subseteq Z(B_F).
\]
\end{enumerate}
The cleaned statement focuses on the third assertion and asks additionally whether the corresponding multivariable monodromy conjecture holds. If the intended monodromy assertion is support-locus containment, its formulation is
\[
 \operatorname{Exp}(H)\subseteq V_F
 \quad\text{for every such polar hyperplane }H.
\]
This interpretation is conditional: the supplied record does not specify the multivariable motivic zeta function, its local or nearby-point convention, the coefficient ring and meaning of a pole, or the precise corresponding monodromy conjecture. In particular, actual polar hyperplanes must not silently be replaced by candidate hyperplanes from a resolution. Thus the displayed pole assertions are formulation schemas pending these source definitions.

The supplied dated curation labels the problem partially solved and attributes the exponentiated zero-locus/support-locus equality to Budur--van der Veer--Wu--Zhou, while describing the ideal-generation assertion and the associated multivariable motivic-monodromy conjecture as unresolved in general. This is a report of the supplied assessment, not an independently verified theorem statement or current status determination.
\subsection{Short English statement}
For a tuple of holomorphic germs, the source asks whether its Bernstein--Sato ideal has generators built from rational affine forms, whether its exponentiated zero locus equals an appropriate local-system cohomology-support locus, and whether every motivic polar hyperplane lies in that zero locus. The cleaned target emphasizes motivic poles and monodromy. The nonzero-component hypothesis, support-locus convention, and precise meaning of motivic poles and multivariable monodromy require source clarification. Supplied dated curation reports partial progress; this draft requires human mathematical review.
\subsection{Sources}
\begin{itemize}
\item Supplied local UnsolvedMath registry record, permanent ID \(30003228\), title \emph{Multivariable Motivic Poles and Monodromy}, source problem code \texttt{catalogue entry}, in the Oberwolfach Reports Open Problems collection. No individual proposer is named; the proposed year is supplied as 2016.
\item Supplied source citation: \emph{Singularities (2017)}, Oberwolfach Reports, pp.\ 2625--2691, DOI \texttt{10.4171/owr/2016/46}. \url{https://doi.org/10.4171/owr/2016/46}. The literature-source list also identifies this reference as Oberwolfach Report 46. The supplied title year 2017 and proposed/report identifier year 2016 are retained without reconciliation.
\item Budur--van der Veer--Wu--Zhou, \emph{Zero loci of Bernstein-Sato ideals}. Supplied literature reference: \url{https://doi.org/10.1007/s00222-020-01025-x}.
\item Budur, \emph{Bernstein-Sato ideals and local systems}. Supplied literature reference: \url{https://doi.org/10.5802/aif.2939}.
\item Literature assessment date: \texttt{2026-08-21}. This is supplied dated curation, not independently verified. It attributes the exponentiated zero-locus/support-locus equality to Budur--van der Veer--Wu--Zhou and describes the stronger ideal-generation assertion and associated multivariable motivic-monodromy conjecture as unresolved in general.
\item The registry supplies the status labels \texttt{partially\_solved} and \texttt{corrected\_verified}, together with a claim of checking the statement against the official report. These are supplied provenance claims; no source was independently visited or verified for this draft.
\end{itemize}
\end{document}
